\section*{第 \ref{ch:g1:day-and-night} 章——昼夜：太阳}

\begin{solution}{exo:g1:day-and-night:1}
\omterm{def:g1:day-and-night:daynight}{日出}、中午、下午、\omterm{def:g1:day-and-night:daynight}{日落}、夜晚。
\end{solution}

\begin{solution}{exo:g1:day-and-night:2}
中午时太阳最高——此时
\omterm{def:g1:light-and-shadows:shadow}{影子}最短。
\end{solution}

\begin{solution}{exo:g1:day-and-night:3}
总是在同一侧。想一想：明天早上留意
第一缕阳光从哪扇窗户照进来——
它会和今天是同一扇窗户。
\end{solution}

\begin{solution}{exo:g1:day-and-night:4}
这种说法不对。太阳整夜都在照耀——
只是照在地球的另一侧。夜晚来临，是因为旋转的地球上我们这一侧
背离了光线，就像你的后脑勺背对着
灯一样。
\end{solution}

\begin{solution}{exo:g1:day-and-night:5}
对他们来说是夜晚，所以他们很可能正在睡觉——
他们那一侧地球背对太阳，而我们这一侧正面向太阳。
\end{solution}

\begin{solution}{exo:g1:day-and-night:6}
光先照到右脸颊，地球仍在转动：接下来光线
照到整张脸。对你的鼻子而言，此时就是早晨\omterm{def:g1:day-and-night:daynight}{日出}。
\end{solution}

\begin{solution}{exo:g1:day-and-night:7}
是我们脚下的大地。地球自转，带着我们一起转动；
太阳只是\emph{看起来}横穿天空，那是因为我们在转动。
\end{solution}

\begin{solution}{exo:g1:day-and-night:8}
一周中有七次\omterm{def:g1:day-and-night:daynight}{日出}和七个夜晚——
每天各对应一次。
\end{solution}

\begin{solution}{exo:g1:day-and-night:9}
她说得可能有道理：如果我们这一侧一直面向太阳，
这里就会一直是白天。但远处那一侧的孩子
就会一直是黑夜——他们不会感谢佐伊的这个想法！
\end{solution}
